\chapter{Continuous Scaling and the Arithmetic Zero Operator}
\label{v4-arith:w9-r2:archimedean-wilson-surface}

Continuous Mellin scaling has no eigenvectors. Occupation energy has
a complete eigenbasis. This spectral distinction prevents their
unitary intertwining. A separate diagonal operator on the zero
occurrences gives an exact arithmetic exponential and its purity
criterion, with no identification of these different carriers.

\section{Continuous Mellin scaling}
\begin{definition}[scaling and its domain]
\label{v4-arith:w9-r2:def:scaling}
Set $H_c=L^2(\mathbb R_{>0},dx/x)$ and
$(U_a f)(x)=f(e^{-a}x)$ for $a\in\mathbb R$.
In the coordinate $t=\log x$, define
\[
 D_c=-i\partial_t,\qquad
 \operatorname{Dom}D_c=\{f:f(e^t)\in H^1(\mathbb R,dt)\}.
\]
Then $U_a=e^{-iaD_c}$.
The derivative convention is $D_c=i\,\partial_a U_a|_{a=0}$.
\end{definition}

\begin{theorem}[Mellin unitarity]
\label{v4-arith:w9-r2:lem:mellin}
The transform
\[
 (\mathcal M_-f)(u)=\frac1{\sqrt{2\pi}}
                  \int_0^\infty f(x)x^{-iu}\frac{dx}{x}
\]
extends to a unitary $H_c\to L^2(\mathbb R,du)$ and sends $D_c$
to multiplication by $u$ on its maximal domain.
It has no eigenvectors. The normal operator
$\Theta_c=1/2+iD_c$ has continuous spectrum $1/2+i\mathbb R$.
\end{theorem}
\begin{proof}
The logarithmic coordinate is unitary from $H_c$ to $L^2(dt)$.
Plancherel and integration by parts give the transform and generator.
The weak $H^1$ domain maps exactly to $\{v:uv\in L^2(du)\}$.
An eigenvector of multiplication by $u$ would be supported on a
singleton, which has Lebesgue measure zero. The same argument applies
to $1/2+iu$.
\end{proof}

\begin{definition}[continuous scaling exponential]
\label{v4-arith:w9-r2:def:F-infty}
On the continuous scaling carrier, put
\[
 F_\infty:=e^{-\Theta_c}=e^{-1/2}U_1.
\]
This bounded normal operator satisfies
$|F_\infty|=e^{-1/2}1$ unconditionally.
\end{definition}

\section{The occupation and primary-sector comparisons}
\begin{theorem}[obstruction to a unitary energy comparison]
\label{v4-arith:w9-r2:lem:L0-Mellin}
Let $H_d$ have a complete orthonormal eigenbasis for an operator
$L_d$. There is no unitary $V:H_d\to H_c$ with
$VL_dV^{-1}=\Theta_c$, including equality of domains.
In particular, the ordinary Fock energy and the integer occupation
operator cannot satisfy this identity.
\end{theorem}
\begin{proof}
If $L_d e=\lambda e$ with $\|e\|=1$, the claimed identity gives
$\Theta_c Ve=\lambda Ve$. This contradicts the absence of
continuous-scaling eigenvectors. The vacuum alone is sufficient.
\end{proof}

The constructed compact-test map to the zero carrier is
$E:\mathscr A\to K_Z$ of Theorem~\ref{v4-arith:sc:zero-representation}.
Its exact comparison is signed. A Petersson norm of an actual cusp
form is the different form computed by the Rankin--Selberg integral.
Neither of these statements supplies a map from an arithmetic
primary sector to $H_c$.

\begin{theorem}[the correctly typed conjugacy formula]
\label{v4-arith:w9-r2:thm:conjugacy}
Suppose $H_a$ is a Hilbert space, $L_a$ a closed normal operator,
and $V:H_a\to H_c$ a unitary with
$V\operatorname{Dom}L_a=\operatorname{Dom}\Theta_c$ and
$VL_a=\Theta_c V$ on that domain.
Then
\[
 e^{-L_a}=V^{-1}F_\infty V.
\]
The exponential on the left is bounded, and its absolute value is
$e^{-1/2}1$. The preceding obstruction excludes this hypothesis when
$L_a$ is the ordinary Fock or occupation energy.
\end{theorem}
\begin{proof}
Unitary conjugacy preserves the spectral projections of normal
operators. Applying the bounded function $e^{-s}$ on the common
vertical-line spectrum proves the formula and absolute-value claim.
\end{proof}

\section{The zero exponential and its actual spectrum}
\begin{definition}[arithmetic zero exponential]
For the diagonal operator $\Theta_Z$ of
Theorem~\ref{v4-arith:sc:normalized-det}, put
\[
 F_Z=e^{-\Theta_Z}=\operatorname{diag}(e^{-\rho})
 \quad\hbox{on }K_Z=\ell^2(Z).
\]
Because $0\le\Re\rho\le1$, this operator is bounded and invertible,
with $\|F_Z\|\le1$ and $\|F_Z^{-1}\|\le e$.
\end{definition}

\begin{proposition}[spectral measure and failure of local finiteness]
\label{v4-arith:w9-r2:thm:spectral-identification}
For a Borel set $B\subset\mathbb C$, the spectral projection of
$F_Z$ is
\[
 (P_Z(B)u)_\rho=\mathbf1_B(e^{-\rho})u_\rho.
\]
Its spectrum is the closure of $\{e^{-\rho}:\rho\in Z\}$.
The counting pushforward $\sum_\rho\delta_{e^{-\rho}}$ is not
locally finite on $\mathbb C$. Thus it is not an entire-function
zero divisor.
For any fixed $s$, the product
$\prod_\rho(1-e^{-s}e^{-\rho})$ fails the necessary factor-to-one
condition for an ordinary convergent nonzero product.
\end{proposition}
\begin{proof}
Coordinate functional calculus gives the spectral projection.
The spectrum of a bounded diagonal operator is the closure of its
entries: outside that closure, the inverse diagonal resolvent is
bounded; at a limit point, basis vectors give approximate eigenvectors.
All infinitely many occurrences lie in the compact annulus
$e^{-1}\le|z|\le1$. A locally finite counting measure would have
finite mass on that annulus, a contradiction.
Finally $|e^{-s}e^{-\rho}|\ge e^{-\Re s-1}>0$, so the factors do
not tend to one. The normalized additive determinant
$\Delta_Z(s)=\xi(s)$ remains well-defined by its genus-one factors.
\end{proof}

\begin{theorem}[arithmetic purity and Weil positivity]
\label{v4-arith:w9-r2:thm:purity-ATP-RH}
The following are equivalent:
\[
 |F_Z|=e^{-1/2}1,
 \qquad \Re\rho=1/2\ \hbox{for every zero }\rho,
 \qquad W(f*f^*)\ge0\ \hbox{for every }f\in\mathscr A.
\]
This is a theorem about $F_Z$ on $K_Z$.
The equality $|F_\infty|=e^{-1/2}1$ on $H_c$ is unconditional and
carries no zero-location information without a further comparison.
\end{theorem}
\begin{proof}
On the coordinate vector at $\rho$, $|F_Z|$ has eigenvalue
$e^{-\Re\rho}$. This proves the first equivalence.
The second is Weil's positivity criterion on the exact compact
convolution test algebra.
\end{proof}

\section{A surface observable and the required operator comparison}
\begin{definition}[an oriented disk and an abelian curvature observable]
\label{v4-arith:w9-r2:def:Sigma-infty}
Let $\Sigma_\infty=D^2=\{z\in\mathbb C:|z|\le1\}$, oriented by
$dx\wedge dy$, and give it the two-form
\[
 \omega_\infty=\frac{i}{2\pi}
              \frac{dz\wedge d\bar z}{(1+|z|^2)^2}.
\]
Let $a$ be a smooth real connection one-form on the trivial $U(1)$
bundle over the disk, and let $q\in\mathbb Z$.
The curvature observable is
$\mathcal O_q(a)=\exp(iq\int_{D^2}da)$.
\end{definition}

\begin{proposition}[disk normalization and Stokes comparison]
The disk has $\int_{D^2}\omega_\infty=1/2$, and
$\mathcal O_q(a)=\exp(iq\int_{\partial D^2}a)$.
For a fixed connection, this scalar observable acting on $H_c$
cannot equal the nonscalar operator $F_\infty$.
\end{proposition}
\begin{proof}
In polar coordinates the area integral is
$2\int_0^1 r(1+r^2)^{-2}dr=1/2$.
Stokes's theorem gives the boundary expression with its induced
counterclockwise orientation. Gauge transformations on the disk
preserve the curvature. The integral charge makes the boundary
holonomy independent of the angular representative.
Choose a nonzero logarithmic test supported in an interval of length
less than one. Its translate by one has disjoint support, so $U_1$
and hence $F_\infty=e^{-1/2}U_1$ is not scalar.
\end{proof}

\begin{definition}[a specified quantum surface operator]
\label{v4-arith:w9-r2:def:F-infty-WS}
A quantum surface comparison consists of a Hilbert space $H_\Sigma$,
a bounded operator $F_\infty^{\mathrm{WS}}$ on it, and a unitary
$V_\Sigma:H_\Sigma\to H_c$ satisfying
$V_\Sigma F_\infty^{\mathrm{WS}}=F_\infty V_\Sigma$.
These are additional data. The scalar curvature observable alone
does not construct them.
\end{definition}

\begin{proposition}[transport of a specified surface operator]
\label{v4-arith:w9-r2:thm:WS-equals-scaling}
For the specified quantum surface comparison,
$F_\infty^{\mathrm{WS}}=V_\Sigma^{-1}F_\infty V_\Sigma$ and
$|F_\infty^{\mathrm{WS}}|=e^{-1/2}1$.
A further arithmetic comparison with $F_Z$ is required before its
purity can be identified with Weil positivity.
\end{proposition}
\begin{proof}
Multiply the defining intertwining identity by $V_\Sigma^{-1}$.
Unitary conjugacy preserves the absolute value.
The arithmetic criterion concerns the separate coordinate operator
$F_Z$ and follows from the preceding purity theorem.
\end{proof}
